\section{Conclusion}\label{sec:conclusion}

This paper presented Beacon, a direct multi-horizon gradient-boosting forecaster for
day-ahead hourly load in the NEMA zone of ISO New England. Rather than rolling a
single near-term model forward, Beacon trains $24$ independent CatBoost models, one per
horizon, each reading a $168$-hour lookback window together with the calendar and
\emph{forecasted weather at its own target hour}. This design removes recursive error
compounding and makes tomorrow's forecasted temperature --- the dominant driver of
day-ahead load --- directly available to every horizon model.

The contributions are fourfold: a direct per-horizon design that lowers average MAE by
$80.6\%$ relative to a single-model roll-out; target-hour exogenous features as the
lever that cuts day-ahead MAE from $183$ to $77$~MW; identification and correction of a
train/serve weather-source mismatch that had nearly doubled live error; and a
horizon-matched, leakage-audited evaluation protocol against an operational ISO
forecast. On a full year of held-out 2025 data, Beacon attains a day-ahead (h=24) MAE
of $76.7$~MW. In a live, horizon-matched comparison over a $30$-day window, it reaches
a day-ahead MAE of $89.3$~MW against ISO-NE's $92.6$~MW, slightly beating the
operational forecast on MAE and MAPE while remaining marginally behind on $R^2$. A
shuffled-target test confirms the models exploit genuine feature--target structure
rather than leakage.

That a keyless, fully reproducible pipeline built on free public data can match and
modestly beat a professional operational forecast --- once the comparison is made
honestly --- suggests that careful feature construction, source discipline, and honest
evaluation matter as much as model complexity. We hope the methodological lessons,
particularly the two pitfalls this work highlights, prove useful to practitioners
deploying weather-informed load forecasters.
